A **2-coloring** is a map $c:\mathbb R^2\to\{\text{red},\text{blue}\}$.

Two triangles $T,T'$ in $\mathbb R^2$ are **congruent** if there is an isometry (composition of translation, rotation, reflection) mapping $T$ to $T'$.

Given a triangle $T$, say that $T$ is **Ramsey for the coloring $c$** if there exists a monochromatic congruent copy of $T$, i.e. an isometric image $T'$ of $T$ whose three vertices all have the same color.

The statement to prove can be phrased as:

> For every 2-coloring $c$ of $\mathbb R^2$, the set of congruence classes of triangles $T$ that are *not* Ramsey for $c$ has size at most $1$.

The problem text also notes there are colorings (“alternating strips”) where one equilateral triangle must be excluded, and that Shader proved the statement for a fixed right triangle.

---
